\section{引言}\label{sec:intro}

BIKE 等 QC-MDPC 方案建立在迭代式比特翻转（BF）译码之上，其罕见的失败塑造着它们的安全论证~\cite{misoczki2013mdpc,bike2024spec,alagic2025nist}。Level-1 的目标是解码失败率（DFR）约 $2^{-128}$，远低于直接仿真所能触及的水平~\cite{sendrier2019dfr,vasseur2021thesis}。失败并不均匀铺开：它们聚集在特定密钥上~\cite{guo2016key,wang2023weak,nosouhi2023weakkey}，而能检测或计时这些失败的攻击者可恢复密钥~\cite{eaton2018timing}。译码器逐轮的阈值构成一套调度，调整它无须触碰密钥生成。因此我们检验一个带条件的问题：对给定密钥，刻意改变调度是否会改变失败风险。

在刻画 Vasseur 参考实现中的 CLASSIC BF 译码器~\cite{vasseur2021thesis}（记作 C0）时，我们注意到它最早的那次决策异常敏感。首轮作用于未经触碰的综合征，因此首轮的任何失误都会进入此后每一轮的状态。把首轮阈值加 2，条件失败率在随机条件下降低 4.15 倍，在对抗探针下降低 5.47 倍；把阈值减 1，失败率反而上升约八倍。两单位偏移若从第 2 轮起持续施加则具破坏性，只限于一轮时却有益。

我们围绕配对实例重建实验。一个实例是一把密钥 $H$、一个错误 $e$ 及其综合征 $s=He^{\top}$；每个调度变体都译码同一个三元组。FRO-$k$ 只在第 1 轮把阈值抬高 $k$，FRONEG 则把阈值减 1。比较用风险比（RR）汇总，即基线失败率除以该变体的失败率，取值大于 1 有利于该变体。由于失败聚集在密钥上，区间与检验重抽样密钥而非错误。引用的所有比率都来自失败频繁到足以计数的区间；我们对生产尾部不作主张。

本研究有四项贡献。第一，首轮是一扇双向窗口。保守偏移 $+1$ 至 $+3$ 使失败率下降 2.3--5.5 倍，而宽松侧在两个计数内崩塌（第~\ref{sec:formal} 节）。第二，同一偏移的符号随施加轮次与持续方式改变。$+2$ 只在第 2 轮施加一次时带来小收益，从第 2 轮起持续施加则给出 RR 0.0168（第~\ref{sec:temporal} 节）。第三，机制是不对称的。FRO-2 把错误的首轮翻转降低约十倍，而正确翻转降得少得多（第~\ref{sec:mechanism} 节）。第四，只要把错误重量抬进可测区间，窗口就在生产块尺寸上复现；在生产重量本身，只有决策余量与收敛代价可测（第~\ref{sec:production} 节）。最后是一个设计好的负结果。一个首轮代理模型未能预测另一套调度的效应（第~\ref{sec:counterfactual} 节）。

第~\ref{sec:background} 节搭建 QC-MDPC BF 译码、相关工作与首轮窗口的一阶模型。第~\ref{sec:method} 节给出配对设计，第~\ref{sec:results} 节给出结果。第~\ref{sec:discussion} 节解读结果并陈述局限；第~\ref{sec:conclusion} 节收束。
